\section{4. Beginning the adventure}

\subsection{Introducing the dungeon}

Present the dungeon as an ancient magical construction, divine trial, abandoned experiment or creation of an eccentric archmage. Its origin need not be explained. The characters should enter voluntarily, knowing that valuable treasure lies somewhere within. Their motivation may be wealth, a legendary artifact, a commission from a patron or recovery of something important to their campaign.

\begin{quote}\itshape Beyond the stone archway lies a perfectly square chamber. Four doors stand at equal distances from one another, each set into the center of a wall. There are no windows, no visible passages beyond the doors, and no indication of how far beneath the earth you have travelled.\end{quote}

\begin{quote}\itshape The air is still, yet somewhere beyond the walls you hear the distant grinding of stone against stone.\end{quote}

\begin{quote}\itshape Whatever this place is, it does not feel abandoned. It feels as though it has been waiting.\end{quote}

\subsection{The entrance}

The entrance is the only permanently stationary room. It is safe and cannot collapse. Its four door pairs begin unlocked and untrapped, making initial exploration accessible. As surrounding chambers shift, the opposite faces may become locked or trapped. Returning is therefore not necessarily as straightforward as leaving. The entrance serves as the party's geographical anchor and eventual escape point.
